\section{Numerics}
\label{SEC:Numerics}

%%%%%%%%%%%%%%%%%%%%%%%%%%%%%%%%%%%%%%%%%%%%%%%%%%%%%%%%%%%%%%%%%%%%%%%%%%%%%%%%

Nearly all numerical methods with greater than first-order accuracy
will exhibit some sort of numerical instability when used to solve
nonlinear differential equations such as the compressible
Navier-Stokes equations.
%
This is especially true for high-order methods
(i.e., greater than second-order) which can generate errors
when any of the modes arising from the nonlinear terms are
not adequately represented by the numerical scheme.
%
Within a given element, this representation is defined by
the stencil size for finite-difference methods or
the set of basis functions (\thatis polynomial order or modes)
for DG or FR methods.

%%%%%%%%%%%%%%%%%%%%%%%%%%%%%%%%%%%%%%%%%%%%%%%%%%%%%%%%%%%%%%%%%%%%%%%%%%%%%%%%

In these cases, the energy from the unresolved modes are aliased
onto the lower modes producing aliasing errors.
%
The effect these errors can have on a given simulation is
highly unpredictable, and many researchers have shown aliasing errors
can produce severe instabilities within a problem\cite{Kirby_2003, Kirby_2006,
Kopriva_2006, Hesthaven_2010, Castonguay_2011, Gassner_2013, Karniadakis_2013,
Williams_2013, Witherden_2014}.
%\cite{21–23, 25–31}%% From AIAA-2022-1813
%
Preventing aliasing and other numerical instabilities is usually
accomplished through the addition of some type of numerical dissipation.

%%%%%%%%%%%%%%%%%%%%%%%%%%%%%%%%%%%%%%%%%%%%%%%%%%%%%%%%%%%%%%%%%%%%%%%%%%%%%%%%

By default, all filtering\slash limiting methods that are enabled for a
simulation are applied just after updating the solution at the end of each
Runge-Kutta stage.
%
The sponge filter, global filter, and local limiter\slash filter also provide
the following respective inputs to alter the frequency in which they are
applied:
\inp[nmlLimiter::SpongeFilterInterval]{Sponge\_Filter\_Interval},
\inp[nmlLimiter::GlobalFilterInterval]{Global\_Filter\_Interval}, and
\inp[nmlLimiter::LocalLimiterInterval]{Local\_Limiter\_Interval}.
%
However, the application frequency for the positivity-preserving limiter,
entropy limiter, and Thermofix cannot be adjusted so they will always be applied
at the end of each Runge-Kutta stage.
%
When it comes time to apply these methods to the newly updated solution, they
are applied in the following order:
%
\begin{enumerate}[parsep=0pt,itemsep=1pt]
\item
  Sponge filter
\item
  Global filter (if $\inptext{Global\_Filter\_Alpha}\leq 1$)
\item
  Positivity-preserving limiter
\item
  Entropy limiter
\item
  Local limiter\slash filter
\item
  Global filter (if $\inptext{Global\_Filter\_Alpha}>1$)
\item
  Thermofix
\end{enumerate}

%%%%%%%%%%%%%%%%%%%%%%%%%%%%%%%%%%%%%%%%%%%%%%%%%%%%%%%%%%%%%%%%%%%%%%%%%%%%%%%%

\noindent
In the following sections, we discuss several important topics related to
numerical stability:

%%%%%%%%%%%%%%%%%%%%%%%%%%%%%%%%%%%%%%%%%%%%%%%%%%%%%%%%%%%%%%%%%%%%%%%%%%%%%%%%

%!While high-order numerical methods provide more accurate simulations,
%!they also contain less inherent numerical dissipation.
%!%
%!This can make them more prone to numerical instability,
%!especially in flows with steep gradients such as shocks.
%!%
%!Here the polynomial representation within an element is unable to
%!capture such discontinuities, with the highest polynomial modes
%!creating under- and over-shoots.
%!%
%!In the following sections, we discuss important topics related to
%!numerical stability:

%%%%%%%%%%%%%%%%%%%%%%%%%%%%%%%%%%%%%%%%%%%%%%%%%%%%%%%%%%%%%%%%%%%%%%%%%%%%%%%%

\begin{itemize}
\item \nameref{SEC:Aliasing}
\item \nameref{SEC:Filtering}
\HIDE{\item \nameref{SEC:Limiting}}
\item \nameref{SEC:Markers}
\end{itemize}

%%%%%%%%%%%%%%%%%%%%%%%%%%%%%%%%%%%%%%%%%%%%%%%%%%%%%%%%%%%%%%%%%%%%%%%%%%%%%%%%

\subsection{Aliasing Error}
\label{SEC:Aliasing}

%%%%%%%%%%%%%%%%%%%%%%%%%%%%%%%%%%%%%%%%%%%%%%%%%%%%%%%%%%%%%%%%%%%%%%%%%%%%%%%%

\key*{Aliasing error}[aliasing error] occurs when the polynomial basis
functions are insufficient to represent the nonlinear terms.
%
Consider the simplified example shown in \fref{FIG:Aliasing},
in which the exact solution $u$ is given by a fourth order polynomial
and the exact flux $u^2$ is given by an eighth order polynomial.
%
In this case, a fourth order polynomial for the numerical solution
can be used to exactly represent the physical solution.
%
Evaluating the flux at the solution points will also result in a
fourth order polynomial that matches the exact flux at those points.
%
However, at any other location, the fourth order polynomial flux
is unable to accurately approximate the eighth order exact flux.
%
The area between the exact flux and polynomial flux curves represents
the aliasing error.

%%%%%%%%%%%%%%%%%%%%%%%%%%%%%%%%%%%%%%%%%%%%%%%%%%%%%%%%%%%%%%%%%%%%%%%%%%%%%%%%

\begin{figure}[htb]
  \begin{centering}
    \insertimage[width=4.0in]{Aliasing/Aliasing}
    \caption{Example of aliasing error.}
    \label{FIG:Aliasing}
  \end{centering}
\end{figure}

%%%%%%%%%%%%%%%%%%%%%%%%%%%%%%%%%%%%%%%%%%%%%%%%%%%%%%%%%%%%%%%%%%%%%%%%%%%%%%%%

\Linecite{Spiegel_2015b} %AIAA-2015-2744
explored one means of reducing aliasing error
using an early version of GFR.
%
Instead of directly computing the flux using the same polynomial
order and Gauss points as the solution, the flux is evaluated
using a higher quadrature order
(\thatis a higher order polynomial and its corresponding Gauss points).
%
This technique of over-integration by over-sampling at the quarature points
requires that the solution polynomial for the conserved variables
be evaluated (\thatis interpolated) at the additional set of quadrature points
in order to use those values in computing the flux.
%
Then that higher order quadrature polynomial is projected
(in a least squares sense) onto a flux polynomial of the same order
as the solution.
%
This reduces the {\em overall error} between the exact flux and
polynomial flux curves, but in doing so the {\em specific error}
between the exact flux and polynomial flux curves evaluated at the
solution Gauss points will now be non-zero.

%%%%%%%%%%%%%%%%%%%%%%%%%%%%%%%%%%%%%%%%%%%%%%%%%%%%%%%%%%%%%%%%%%%%%%%%%%%%%%%%

The main disadvantage to reducing aliasing error in this manner is
the additional cost involved in supersampling at the additional
\key{quadrature points}[points!quadrature] and then projecting that
polynomial down to the solution order.
%
For that reason, subsequent development of GFR (\thatis the viscous terms)
has not incorporated the quadrature evaluation.

%%%%%%%%%%%%%%%%%%%%%%%%%%%%%%%%%%%%%%%%%%%%%%%%%%%%%%%%%%%%%%%%%%%%%%%%%%%%%%%%



\subsection{Filtering}
\label{SEC:Filtering}
\key*{}[filtering]

%%%%%%%%%%%%%%%%%%%%%%%%%%%%%%%%%%%%%%%%%%%%%%%%%%%%%%%%%%%%%%%%%%%%%%%%%%%%%%%%

In spectral-element methods, the solution within an individual element
can be transformed into a hierarchical set of modal coefficients which
are representative of the energy contained in each mode.
%
A well-resolved numerical solution will have minimal energy contained
within the higher modes, and the modal coefficients should have an
exponential decay rate from lowest to highest modes.
%\cite{21, 22, 32, 33}%% From AIAA-2022-1813
%
For this reason, an exponential filter is used as a source of
numerical dissipation for stabilizing the simulations.

%%%%%%%%%%%%%%%%%%%%%%%%%%%%%%%%%%%%%%%%%%%%%%%%%%%%%%%%%%%%%%%%%%%%%%%%%%%%%%%%
%% Note for future reference:
%%
%% For the diagonal filter matrix sigma(m) in GFR,
%% see: Spiegel AIAA-2022-1813 Eq.5,
%% but that index goes from m=0:N_M "for the N_M polynomial modes (p.2)".
%% The polynomial order is degree P (p.2) and 
%% the number of solution points is N_P=P+1.
%% His previous paper numbers the modal coefs 0:P and the nodal coefs 1:P=1.
%% Since a polynomial of degree P has terms x^0 + x^1 + ... x^P
%% it will have just as many modes 0:P with 0 being the constant mode.
%% In this context, it seems N_M=P even though it is not explicitly defined.
%% However, in practice the arrays in GFR begin at index 1.
%% 
%% Hesthaven & Warburton (2010) \cite{Hesthaven_2010}
%% p.20:
%%   N=polynomial order; N_P=polynomial points; so N_p=N+1.
%%   The modal coefs are 1:N_P and the nodal coefs are 1:N_P too.
%% p.130:
%%   Defines sigma(eta) with eta=(n-1)/N; n=1...N_p;
%%   Their Matlab code on that page also provides a
%%     1-based diagonal filter matrix.
%%
%% Sheshadri (2016) PhD \cite{sheshadri16}
%% p. 11:  The nodal coefs are n=0:P.
%% p. 99:  The modal coefs are n=0:P too.
%% p.107:
%%   Defines sigma(eta) with eta=n/P; n=0...P;
%%
%% We use Hesthaven's notation here, because it is consistent with the
%% 1:P+1 modal form of the polynomial used in our introduction.

This diagonal exponential filter matrix is given by\cite{Hesthaven_2010}
\begin{align}
    \label{EQN:ExpFilterMatrix}
    \sigma (n)
 =& \begin{cases}
       1
    &  0 \leq \eta \leq \eta_c
    \\
       \exp \left[ -\alpha \left( \frac{\eta-\eta_c}{1-\eta_c} \right)^s
            \right]
    &  \eta_c \leq \eta \leq 1
    \end{cases}
  &
    \begin{aligned}
       \eta   = & \;\text{modal fraction}, (n-1) / \PO
    \\ \eta_c = & \;\text{modal fraction cutoff}
    \\ \alpha = & \;\text{filter strength}
    \\ s      = & \;\text{filter order}
    \end{aligned}
\end{align}
%
where $n=1,\ldots,\PO+1$.

%%%%%%%%%%%%%%%%%%%%%%%%%%%%%%%%%%%%%%%%%%%%%%%%%%%%%%%%%%%%%%%%%%%%%%%%%%%%%%%%

%Three user-specified parameters control the aggressiveness of the filter,
%and their effects are illustrated in \fref{FIG:Filter}.
%%
%%
%(1) The parameter $\eta_c$ specifies the modal fraction cutoff above which
%filtering will be applied.
%%
%Everything below that modal fraction cutoff will be entirely preserved.
%%
%Most importantly, the lowest mode
%(which defines the average value in the cell)
%remains unchanged in order to ensure conservation of
%mass, momentum, and energy.
%%
%%
%(2) The parameter $\alpha$ specifies the strength of the filter.
%As $\alpha$ increases, the filter value at the highest mode(s)
%becomes smaller, which means that more damping is applied.
%%
%%
%(3) Lastly, the parameter $s$ specifies the filter order
%which controls the exponential decay rate
%(\thatis how parabolic the filter function is on a log plot).
%As $s$ increases, filtering has substantially less of an effect on
%the lower modes.

%%%%%%%%%%%%%%%%%%%%%%%%%%%%%%%%%%%%%%%%%%%%%%%%%%%%%%%%%%%%%%%%%%%%%%%%%%%%%%%%

Three user-specified parameters, $[\eta_c, \alpha, s]$, control the
aggressiveness of the filter, and their effects are illustrated in
\fref{FIG:Filter}.\\[12pt]

%%%%%%%%%%%%%%%%%%%%%%%%%%%%%%%%%%%%%%%%%%%%%%%%%%%%%%%%%%%%%%%%%%%%%%%%%%%%%%%%

\begin{figure}[htb]
  \begin{centering}
    \begin{subfigure}{2.0in}
      \insertimage[width=\textwidth]{Filter/Filter_etac}
      \caption{Effect of $\eta_c$}
      \label{FIG:Filter_etac}
    \end{subfigure}%
    \begin{subfigure}{2.0in}
      \insertimage[width=\textwidth]{Filter/Filter_alpha}
      \caption{Effect of $\alpha$}
      \label{FIG:Filter_alpha}
    \end{subfigure}%
    \begin{subfigure}{2.0in}
      \insertimage[width=\textwidth]{Filter/Filter_s}
      \caption{Effect of $s$}
      \label{FIG:Filter_s}
    \end{subfigure}%
    %%%%%%%%%%%%%%%%%%%%%%%%%%%%%%%%%%%%%%%%%%%%%%%%%%%%%
    \caption{Exponential filter function, $\sigma$.}
    \label{FIG:Filter}
  \end{centering}
\end{figure}

%%%%%%%%%%%%%%%%%%%%%%%%%%%%%%%%%%%%%%%%%%%%%%%%%%%%%%%%%%%%%%%%%%%%%%%%%%%%%%%%
\begin{tabular}{%
      @{\hspace{2.0em}} l%
      @{\hspace{1.0ex}} c%
      @{\quad} >{\raggedright\arraybackslash}p{0.8\linewidth}%
      @{}}
  %
  $\eta_c$ & $\coloneq$ &
  This parameter specifies the modal fraction cutoff above which filtering will
  be applied.
  %
  Everything below that modal fraction cutoff will be entirely preserved.
  %
  Most importantly, the lowest mode (which defines the average value in the
  cell) remains unchanged in order to ensure conservation of mass, momentum, and
  energy.
  %
  Valid input values for this parameter are real values on the interval
  $[0.0,1.0]$, although a value of $1.0$ is the same as using no filter at all
  because it will preserve all modes of the solution.
  %
  A value of $0.0$ will apply some amount of filtering to all modes except the
  lowest\slash constant mode.\\
  %
  & & \\
  %   
  $\alpha$ & $\coloneq$ & 
  This parameter specifies the strength of the filter, \thatis the highest mode
  will be multiplied by the value $\exp\left(-\alpha\right)$.
  %
  As $\alpha$ increases, the filter value at the highest mode(s) becomes
  smaller, which means that more damping is applied.
  %
  Valid input values for this parameter are real values between $0.0$ and
  roughly $36.0$, where a value of $36.0$ is nearly equivalent to multiplying
  the highest mode by 64-bit machine epsilon, \thatis
  $\ln\left(2^{-52}\right)\approx-36.0$.\\
  %
  & & \\
  %
  $s$ & $\coloneq$ & 
  This parameter specifies the filter order which controls the exponential decay
  rate (\thatis how parabolic the filter function is on a log plot).
  %
  As $s$ increases, filtering has substantially less of an effect on the lower
  modes.
  %
  Valid input values for this parameter are positive integers greater than or
  equal to 2; however, anything greater than 32 will essentially result in a
  filter that multiplies the highest mode by $\exp\left(-\alpha\right)$ applies
  negligible filtering to any lower modes that are larger than the cutoff mode.
  %
  In practice, commonly used values are $[2,4,8,12,16,32]$.
  %
\end{tabular}
\\[12pt]

%%%%%%%%%%%%%%%%%%%%%%%%%%%%%%%%%%%%%%%%%%%%%%%%%%%%%%%%%%%%%%%%%%%%%%%%%%%%%%%%

GFR uses a three-pass filtering approach.
%
First, a \key*{\emph{sponge filter}}[filtering!sponge filter] can be applied
to the downstream region starting from a user-specified x-location.
This filter region is often used to damp pressure fluctuations so they do not
reflect from the outflow boundary and contaminate the upstream solution.
%
Second, a {\key*{\emph{local filter}}[filtering!local filter] can be applied
to those cells identified by the
\key[key:marker]{troubled cell marker}[marker] in \sref{SEC:Markers}.
%
Usually these cells require some significant level of filtering
due to a shock or some local instability.
%
Third, a weak \key*{\emph{global filter}}[filtering!global filter] can be
applied to all cells in the domain to remove any high-frequency noise in the
solution; NOTE, use of the global filter should be avoided if possible since it
filters \emph{everything} including cells that do not need filtering.

%%%%%%%%%%%%%%%%%%%%%%%%%%%%%%%%%%%%%%%%%%%%%%%%%%%%%%%%%%%%%%%%%%%%%%%%%%%%%%%%

It is recommended to check for troubled cells and perform 
filtering operations after each Runge-Kutta subiteration.
%
This is the highest filtering frequency possible, with other researchers
opting to filter only after the final Runge-Kutta stage or after a
certain number of time steps.
%
Filtering at such a high frequency allows for the use of a larger
time step, but it can reduce the accuracy of the solution.
%
This decrease in accuracy is usually outweighed by the increased
time step size, which is often the most limiting computational factor
in time-to-solution performance.

%%%%%%%%%%%%%%%%%%%%%%%%%%%%%%%%%%%%%%%%%%%%%%%%%%%%%%%%%%%%%%%%%%%%%%%%%%%%%%%%

\SeeAlso{}
\noindent%
\begin{tabular}{@{\qquad} lll @{}}
\begin{varwidth}[t]{0.3\textwidth}
\inp[nmlLimiter::GlobalFilterInterval]{Global\_Filter\_Interval}, \\
\inp[nmlLimiter::GlobalFilterOption]{Global\_Filter\_Option}, \\
\inp[nmlLimiter::GlobalFilterOrder]{Global\_Filter\_Order}, \\
\inp[nmlLimiter::GlobalFilterAlpha]{Global\_Filter\_Alpha}, \\
\inp[nmlLimiter::GlobalFilterEtac]{Global\_Filter\_Etac}, \\
%
\end{varwidth}
&
\begin{varwidth}[t]{0.3\textwidth}
\inp[nmlLimiter::LocalLimiterInterval]{Local\_Limiter\_Interval}, \\
\inp[nmlLimiter::LocalFilterOrder]{Local\_Filter\_Order}, \\
\inp[nmlLimiter::LocalFilterAlpha]{Local\_Filter\_Alpha}, \\
\inp[nmlLimiter::LocalFilterEtac]{Local\_Filter\_Etac}, \\
%
\end{varwidth}
&
\begin{varwidth}[t]{0.3\textwidth}
\inp[nmlLimiter::SpongeFilterInterval]{Sponge\_Filter\_Interval}, \\
\inp[nmlLimiter::SpongeFilterOption]{Sponge\_Filter\_Option}, \\
\inp[nmlLimiter::SpongeFilterOrder]{Sponge\_Filter\_Order}, \\
\inp[nmlLimiter::SpongeFilterAlpha]{Sponge\_Filter\_Alpha}, \\
\inp[nmlLimiter::SpongeFilterEtac]{Sponge\_Filter\_Etac}, \\
\inp[nmlLimiter::SpongeFilterXStart]{Sponge\_Filter\_X\_Start}\end{varwidth}
\end{tabular}

%%%%%%%%%%%%%%%%%%%%%%%%%%%%%%%%%%%%%%%%%%%%%%%%%%%%%%%%%%%%%%%%%%%%%%%%%%%%%%%%



\HIDE{
\subsection{Limiting}
\label{SEC:Limiting}
}

%%%%%%%%%%%%%%%%%%%%%%%%%%%%%%%%%%%%%%%%%%%%%%%%%%%%%%%%%%%%%%%%%%%%%%%%%%%%%%%%



\subsection{Troubled Cell Markers}%
\label{SEC:Markers}%
\key*[key:marker]{}[marker]%

\IDEA{Add Ducros and Ducros+Jump markers.}

%%%%%%%%%%%%%%%%%%%%%%%%%%%%%%%%%%%%%%%%%%%%%%%%%%%%%%%%%%%%%%%%%%%%%%%%%%%%%%%%

%! Blindly applying filtering or limiting to every element in the domain
%! can cause a reduction in numerical accuracy.
%! %
%! Therefore, we wish to restrict such treatments to only those elements
%! that are experiencing numerical instability.
%! %
%! This is done through the use of a troubled cell marker.
%! %
%! A few different types of markers are available for use and are
%! described below.

%%%%%%%%%%%%%%%%%%%%%%%%%%%%%%%%%%%%%%%%%%%%%%%%%%%%%%%%%%%%%%%%%%%%%%%%%%%%%%%%

Instabilities within a given simulation are usually highly-localized
to very small regions of the grid domain.
%
Frequently, just a single cell out of perhaps millions of cells will
entirely contain the instability at a particular time step of a simulation.
%
Indiscriminately applying numerical dissipation to all cells in a
simulation for stabilization purposes will add unwanted dissipation
in places where it is not needed, severely affecting the overall accuracy
of the solution.
%
Therefore, a method is needed to identify only those cells whose
solutions show signs of becoming unstable, so they can be specifically
targeted with additional numerical dissipation.
%
A few different types of troubled cell markers are available for use
and are described below.


%%%%%%%%%%%%%%%%%%%%%%%%%%%%%%%%%%%%%%%%%%%%%%%%%%%%%%%%%%%%%%%%%%%%%%%%%%%%%%%%

\subsubsection{Resolution Marker}
\label{SEC:ResolutionMarker}
\key*{}[marker!resolution marker]

%%%%%%%%%%%%%%%%%%%%%%%%%%%%%%%%%%%%%%%%%%%%%%%%%%%%%%%%%%%%%%%%%%%%%%%%%%%%%%%%

%! The idea behind the resolution marker is that numerical instabilities
%! often appear in the highest polynomial mode(s).

As discussed in \sref{SEC:Filtering}, numerical instabilities often
appear in the highest polynomial mode(s).
%
An easy way to distinguish between well-resolved and possibly unstable
cells is to examine the highest mode of the solution within each cell.
 
%%%%%%%%%%%%%%%%%%%%%%%%%%%%%%%%%%%%%%%%%%%%%%%%%%%%%%%%%%%%%%%%%%%%%%%%%%%%%%%%

Recall that the \key{modal form}[polynomial!modal form] of the
solution polynomial for a given flow variable is given by \eref{EQN:Modal_Form}
\begin{align*}
       u (r) =& \sum_{n=1}^{N_p} \uModal[n] (t) \psi_n (r)
& \psi_n (r) =& r^{n-1}
\end{align*}
where $\uModal[n]$ represents the modal coefficients, with $\uModal[1]$ being
the coefficient for the lowest\slash constant mode and $\uModal[N_p]$ being
the coefficient for the highest\slash degree-\PO{} mode.
%
A diagonal filter matrix can be applied to remove the highest mode of the
solution.
\begin{align}
    \bm{\underline{\uModal[]}}
 =& \bm{\sigma} \bm{\uModal[]}
  &
    \sigma_n
 =& \begin{cases}
      0 & n = N_p \\
      1 & \text{otherwise}
    \end{cases}
\end{align}
 
A comparison of this filtered solution to the actual solution will
provide a test for the possibility of numerical instabilities in
the cell.
%
Namely, in a well-resolved solution,
$\bm{\uModal[]} \simeq \bm{\underline{\uModal[]}}$.
%
Any cells in which this comparison shows a difference that is beyond
a certain threshold are marked as ``troubled'' cells.

%%%%%%%%%%%%%%%%%%%%%%%%%%%%%%%%%%%%%%%%%%%%%%%%%%%%%%%%%%%%%%%%%%%%%%%%%%%%%%%%

To do this, we compute the local resolution marker
as\cite{Persson_2006,Barter_2007,Barter_2008}
\begin{align}
    \label{EQN:ResolutionMarker}
    \marker \left( \bm{\uModal[]} , \bm{\underline{\uModal[]}} \right)
 =& \log_{10}
    \left[
      \frac{
        \left( \bm{\uModal[]} - \bm{\underline{\uModal[]}} \right) \cdot
        \left( \bm{\uModal[]} - \bm{\underline{\uModal[]}} \right)
      }{
               \bm{\uModal[]} \cdot        \bm{\uModal[]}
      }
    \right]
\end{align}
where $\bm{\uModal[]}$ is the user-specified solution variable
to check (density is recommended).
%
If all of the conserved variables are specified, then
\eref{EQN:ResolutionMarker} is computed for each one and the
maximum value is used.
%
Note that the numerator represents the norm of the highest mode
and the denominator represents the norm of all the modes combined.
%
The marker will have values in the range $\marker = [-\infty,0]$ where
the lower value indicates that nothing is contained in the highest mode and 
the upper value indicates that everything is contained in the highest mode.

%%%%%%%%%%%%%%%%%%%%%%%%%%%%%%%%%%%%%%%%%%%%%%%%%%%%%%%%%%%%%%%%%%%%%%%%%%%%%%%%

Next, we compute the local scaled marker as described in \sref{SEC:ScaledMarker}.

%%%%%%%%%%%%%%%%%%%%%%%%%%%%%%%%%%%%%%%%%%%%%%%%%%%%%%%%%%%%%%%%%%%%%%%%%%%%%%%%

%!% AIAA-2022-1813 by Spiegel used:
%!\begin{align}
%!    {\cal S} >& - \left( \psi + 4.25 \log_{10} \left( \PO \right) \right) &
%!\end{align}
%!where $\psi$ is a user-specified constant that will depend on the
%!problem and the polynomial order $\PO$.

%%%%%%%%%%%%%%%%%%%%%%%%%%%%%%%%%%%%%%%%%%%%%%%%%%%%%%%%%%%%%%%%%%%%%%%%%%%%%%%%

\SeeAlso{
\inp[nmlLimiter::MarkerOption]{Marker\_Option},
\inp[nmlLimiter::MarkerVariable]{Marker\_Variable}
}


%%%%%%%%%%%%%%%%%%%%%%%%%%%%%%%%%%%%%%%%%%%%%%%%%%%%%%%%%%%%%%%%%%%%%%%%%%%%%%%%



%%%%%%%%%%%%%%%%%%%%%%%%%%%%%%%%%%%%%%%%%%%%%%%%%%%%%%%%%%%%%%%%%%%%%%%%%%%%%%%%

\subsubsection{Jump Marker}
\label{SEC:JumpMarker}
\key*{}[marker!jump marker]

%%%%%%%%%%%%%%%%%%%%%%%%%%%%%%%%%%%%%%%%%%%%%%%%%%%%%%%%%%%%%%%%%%%%%%%%%%%%%%%%

%% Note:
%% - Krivodonova_2004: Original jump marker paper (as far as I can tell)
%% - Burgess_2011aiaa: Provides the jump marker with    log10 scaling.
%% - Burgess_2011thesis: Burgess's PhD that is similar to Burgess_2011aiaa.
%% - Barter_2007:  Provides the jump marker without log10 scaling
%%                 and coefficients (C1,C2,dPsi) for the scaled marker.
%% - Barter_2008:  Barter's PhD that is similar to Barter_2007.

For flows with shocks or other sudden discontinuities,
it is useful to mark those cells that experience a sudden jump
in pressure or density.
%
The jump marker\cite{Krivodonova_2004, Barter_2007, Barter_2008,
Burgess_2011aiaa, Burgess_2011thesis} is a method for identifying such sudden
changes across adjacent cells.
%
Using the notation in \fref{FIG:Jump_Marker}, the jump indicator
can be defined as:
\begin{align}
  \label{EQN:JumpIndicator1}
    J_\kappa
 =& \frac{1}{\|\partial \kappa\|}
    \int_{\partial \kappa}
       \left|
          \frac{ [[ \psi ]] \cdot \vec{n} }
               { \{ \psi \}               }
       \right|
       ds
  & \begin{aligned}
      [[ \psi ]] =& \psi_1 \vec{n}_1 + \psi_2 \vec{n}_2
                 =  \left( \psi_1 - \psi_2 \right) \vec{n}_1 \\
      \{ \psi \} =& \left( \psi_1 + \psi_2 \right) / 2
    \end{aligned}
 \\
  \label{EQN:JumpIndicator2}
 =& \frac{1}{\|\partial \kappa\|}
    \int_{\partial \kappa}
       \left|
         \frac{ \left( \psi_1 - \psi_2 \right) }
              { \left( \psi_1 + \psi_2 \right) / 2 }
       \right|
       ds
 =  \left[ 0 : 2 \right]
\end{align}
%
Here,
$[[\psi]]$ represents the jump across adjacent face values and
$\{\psi\}$ represents the average of those values.
%
The jump indicator $J_\kappa$ can vary between 0 (no jump)
and 2 (infinite jump).

\begin{figure}
  \begin{centering}
    \insertimage[width=0.5\textwidth]{tikz/Jump_Marker}
   %\input{Figures/Jump_Marker}
    \caption{Jump Marker Notation}
    \label{FIG:Jump_Marker}
  \end{centering}
\end{figure}

The jump marker is obtained by taking the logarithm of the value.
\begin{align}
    {\cal M} \left( J_\kappa \right)
 =& \log_{10} \left[ \max \left( J_\kappa , 2^{-23} \right) \right]
 =  \left[ -6.9 : 0.3 \right]
  \label{EQN:JumpMarker}
\end{align}
The jump marker ${\cal M}$ can vary between -6.9 (no jump)
and 0.3 (infinite jump).
%
The scaled marker is then computed as described in \sref{SEC:ScaledMarker}.

%%%%%%%%%%%%%%%%%%%%%%%%%%%%%%%%%%%%%%%%%%%%%%%%%%%%%%%%%%%%%%%%%%%%%%%%%%%%%%%%

\SeeAlso{
\inp[nmlLimiter::MarkerOption]{Marker\_Option},
\inp[nmlLimiter::MarkerVariable]{Marker\_Variable}
}

%%%%%%%%%%%%%%%%%%%%%%%%%%%%%%%%%%%%%%%%%%%%%%%%%%%%%%%%%%%%%%%%%%%%%%%%%%%%%%%%



%%%%%%%%%%%%%%%%%%%%%%%%%%%%%%%%%%%%%%%%%%%%%%%%%%%%%%%%%%%%%%%%%%%%%%%%%%%%%%%%

\subsubsection{MLP Marker}
\label{SEC:MLPMarker}
\key*{}[marker!MLP marker]

%%%%%%%%%%%%%%%%%%%%%%%%%%%%%%%%%%%%%%%%%%%%%%%%%%%%%%%%%%%%%%%%%%%%%%%%%%%%%%%%

The MLP marker is normally computed as part of the
Multidimensional Limiting Process\cite{Park_2014,Park_2016,Park_2017}.
%
It first identifies troubled cells using the
$\PO[1]$-projected MLP condition, which compares the cell vertex (corner)
values to those of the adjacent cells.
%
Then the MLP extrema detector is applied to identify local smooth extrema
near the vertex points and unmark those cells.

%%%%%%%%%%%%%%%%%%%%%%%%%%%%%%%%%%%%%%%%%%%%%%%%%%%%%%%%%%%%%%%%%%%%%%%%%%%%%%%%

When GFR outputs the MLP marker, the value \marker{} represents the
polynomial order of the limited solution,
with possible real values of [0:1,2,3,...,$\PO$].
%
Values less than $\PO$ indicate that the polynomial in the cell
has been reduced to that value to maintain stability.
%
Values between [0:1] indicate that all of the higher modes have been removed
and the linear mode was limited by that fraction in order to remain stable.

%%%%%%%%%%%%%%%%%%%%%%%%%%%%%%%%%%%%%%%%%%%%%%%%%%%%%%%%%%%%%%%%%%%%%%%%%%%%%%%%

\SeeAlso{
\inp[nmlLimiter::MarkerOption]{Marker\_Option},
\inp[nmlLimiter::MarkerVariable]{Marker\_Variable}
}


%%%%%%%%%%%%%%%%%%%%%%%%%%%%%%%%%%%%%%%%%%%%%%%%%%%%%%%%%%%%%%%%%%%%%%%%%%%%%%%%



%%%%%%%%%%%%%%%%%%%%%%%%%%%%%%%%%%%%%%%%%%%%%%%%%%%%%%%%%%%%%%%%%%%%%%%%%%%%%%%%

\subsubsection{Ducros Marker}
\label{SEC:DucrosMarker}
\key*{}[marker!Ducros Marker]

%%%%%%%%%%%%%%%%%%%%%%%%%%%%%%%%%%%%%%%%%%%%%%%%%%%%%%%%%%%%%%%%%%%%%%%%%%%%%%%%

The Ducros marker uses a divergence based indicator together with the jump
marker to identify troubled cells.
%
First, the jump marker is computed using
\eref{EQN:JumpIndicator2,EQN:JumpMarker} and then the scaled jump marker,
$\scaledmarker_{\textrm{jump}}$, is evaluated using
\eref{EQN:ScaledResolutionMarker}.
%
Next, the divergence indicator is computed using the equation:
%
\begin{equation}
  \scaledmarker_{\textrm{div}} = \left(
                                 \frac{2 \left|\nabla\cdot\vec{V}\right|}{a}
                                 \cdot
                                 \frac{h}{\p[+1]{}}
                                 \right)^4
\end{equation}
%
where $h$ is the maximum length of a cell across all Cartesian coordinate
directions and $a$ is the speed of sound.
%
The value of $\scaledmarker_{\textrm{div}}$ for a cell is taken as the maximum
value from all solution points of the cell.
%
Finally, the scaled Ducros marker is computed as
%
\begin{equation}
  \scaledmarker_{\textrm{Ducros}} = \scaledmarker_{\textrm{jump}} \cdot
                                    \scaledmarker_{\textrm{div}}
\end{equation}

%%%%%%%%%%%%%%%%%%%%%%%%%%%%%%%%%%%%%%%%%%%%%%%%%%%%%%%%%%%%%%%%%%%%%%%%%%%%%%%%

Further information about the Ducros marker can be found in
\linecite{Koeplinger_2026}.

%%%%%%%%%%%%%%%%%%%%%%%%%%%%%%%%%%%%%%%%%%%%%%%%%%%%%%%%%%%%%%%%%%%%%%%%%%%%%%%%

\SeeAlso{
\inp[nmlLimiter::MarkerOption]{Marker\_Option},
\inp[nmlLimiter::MarkerVariable]{Marker\_Variable}
}


%%%%%%%%%%%%%%%%%%%%%%%%%%%%%%%%%%%%%%%%%%%%%%%%%%%%%%%%%%%%%%%%%%%%%%%%%%%%%%%%



%%%%%%%%%%%%%%%%%%%%%%%%%%%%%%%%%%%%%%%%%%%%%%%%%%%%%%%%%%%%%%%%%%%%%%%%%%%%%%%%

\subsubsection{Scaling the Marker Results to Identify Troubled Cells}
\label{SEC:ScaledMarker}
\key*{}[marker!scaled marker]

%%%%%%%%%%%%%%%%%%%%%%%%%%%%%%%%%%%%%%%%%%%%%%%%%%%%%%%%%%%%%%%%%%%%%%%%%%%%%%%%

From the local marker, a local {\em scaled} marker can be computed
\begin{align}
    \label{EQN:ScaledResolutionMarker}
    \scaledmarker
 =& \begin{cases}
      0                                  &
      \phantom{\bigl|} \zeta \phantom{\bigr|} \leq -\pi/2
 \\
      \left( 1 + \sin \zeta \right) / 2  & 
               \bigl|  \zeta          \bigr|  \leq  \pi/2
 \\
      1                                  & 
      \phantom{\bigl|} \zeta \phantom{\bigr|} \geq  \pi/2
    \end{cases}
  & 
    \zeta
 =& \frac{\pi}{2} \, \frac{ \marker - \Psi_0 }{ \Delta \Psi }
\end{align}
where
\begin{align}
    \Psi_0 =& -\left[ C_1 + C_2 \log_{10} (\PO) \right]
  & \Delta \Psi =& 0.5
\end{align}

%%%%%%%%%%%%%%%%%%%%%%%%%%%%%%%%%%%%%%%%%%%%%%%%%%%%%%%%%%%%%%%%%%%%%%%%%%%%%%%%

As shown in \fref{FIG:Scaled_Marker}, the scaled marker varies between
0 (region to retain) and 1 (region to filter).
%
Other works sometimes use this smooth variation to modify the
strength of the dissipation applied.
%
In GFR, this smoothly varying scaled marker can be output for
visualization purposes.
%
However, troubled cells are simply identified by
${\cal S} > \scinum{1.0}{-8}$.
%
This approach is more aggressive at identifying troubled cells.

%%%%%%%%%%%%%%%%%%%%%%%%%%%%%%%%%%%%%%%%%%%%%%%%%%%%%%%%%%%%%%%%%%%%%%%%%%%%%%%%

\begin{figure}
  \begin{centering}
    \insertimage[width=4.0in]{Scaled_Marker/Scaled_Marker}
    \caption{Scaled Marker}
    \label{FIG:Scaled_Marker}
  \end{centering}
\end{figure}

%%%%%%%%%%%%%%%%%%%%%%%%%%%%%%%%%%%%%%%%%%%%%%%%%%%%%%%%%%%%%%%%%%%%%%%%%%%%%%%%

%% Description for the old version, when $(C_1,C_2)$ were still inputs:
%
% The default values of the $(C_1,C_2)$ coefficients are
% (4.00,4.25) when using the resolution marker and
% (2.25,3.00) when using the jump marker.
% %
% GFR allows one to either specify $(C_1,C_2,\Delta\Psi)$
% or to directly specify
% \inp[nmlLimiter::MarkerMin]{Marker\_Min} = $\Psi_{min}$
%                                          = $(\Psi_0-\Delta\Psi)$ and
% \inp[nmlLimiter::MarkerMax]{Marker\_Max} = $\Psi_{max}$
%                                          = $(\Psi_0+\Delta\Psi)$
% in which the variation in the scaled marker $\scaledmarker$ occurs.

The default values of the $(C_1,C_2)$ coefficients are
$(4.0,4.25)$ when using the resolution marker and
$(2.25,3.0)$ when using the jump marker.
%
GFR will use $(C_1,C_2,\Delta\Psi)$ to compute 
$\Psi_{min} = \Psi_0-\Delta\Psi$ and
$\Psi_{max} = \Psi_0+\Delta\Psi$
unless the inputs for
(\inp[nmlLimiter::MarkerMin]{Marker\_Min},%
 \inp[nmlLimiter::MarkerMax]{Marker\_Max})
=
($\Psi_{min},\Psi_{max}$) are used instead.

%%%%%%%%%%%%%%%%%%%%%%%%%%%%%%%%%%%%%%%%%%%%%%%%%%%%%%%%%%%%%%%%%%%%%%%%%%%%%%%%

When using the MLP marker, the scaled marker is simply computed as
a binary flag to indicate a troubled cell, \thatis any cell in which
the multidimensional limiting process would reduce the solution
polynomial order.
\begin{align}
    \scaledmarker
 =& \begin{cases}
       0 & \text{if $\marker = \PO$}
    \\ 1 & \text{otherwise}
    \end{cases}
\end{align}

%%%%%%%%%%%%%%%%%%%%%%%%%%%%%%%%%%%%%%%%%%%%%%%%%%%%%%%%%%%%%%%%%%%%%%%%%%%%%%%%

\SeeAlso{
\inp[nmlLimiter::MarkerDpsi]{Marker\_Dpsi},
\inp[nmlLimiter::MarkerMin]{Marker\_Min},
\inp[nmlLimiter::MarkerMax]{Marker\_Max}
}



%%%%%%%%%%%%%%%%%%%%%%%%%%%%%%%%%%%%%%%%%%%%%%%%%%%%%%%%%%%%%%%%%%%%%%%%%%%%%%%%



